\begin{textsample}{POS Dim 1 – pa – Score 47.00 – pa/pa\_em\_6.txt}  \label{ex:f1_pos_039}
2.1.1 Respeito às Diversas Culturas Amazônicas e suas Inter-Relações no Espaço e no Tempo No \textbf{que} \textbf{diz} respeito às diferentes representações e concepções de Amazônia, é importante se considerar \textbf{que} nesta região existem várias “ Amazônias ”.
Todas elas possuem suas especificidades e \textbf{singularidades}, sejam sociais, econômicas, culturais, territoriais, políticas e ambientais, isto é, refletem a interculturalidade do estado do Pará ( GONÇALVES, 2012 ).
A Amazônia precisa ser analisada e problematizada a partir de \textbf{uma} visão complexa e humanizada, \textbf{que} fuja ao paradigma de análise do pensamento ocidental entre sociedade e natureza, pois não se pode dissociar o universo ecológico do humano e vice-versa.
É importante \textbf{que} se considere o conhecimento dos povos \textbf{que} \textbf{vivem} na floresta e \textbf{que} dela retiram a sua sobrevivência, respeitando suas crenças, símbolos e modos de vida.
Esses povos não são apenas moradores deste espaço, mas são acima de tudo produtores e reprodutores dele ( GONÇALVES, 2012 ).
Os povos da floresta, como os indígenas em sua multiplicidade étnica, os caboclos \textbf{ribeirinhos} ( miscigenados ), os quilombolas descendentes dos escravos, os seringueiros descendentes de imigrantes nordestinos, os garimpeiros, os agroextrativistas, como os castanheiros, os pequenos produtores rurais de subsistência, as mulheres quebradeiras de coco babaçu, entre outras coletividades \textbf{que} apresentam relações heterogêneas com a floresta, com os rios e com seus ecossistemas, precisam ter suas \textbf{territorialidades} contempladas por políticas públicas, sobretudo as políticas públicas educacionais.
Esses povos constroem e materializam no território amazônico diversos modos e meios de vida ; cada \textbf{um} desses modos e meios de vida confere, por sua vez, diferentes usos do território.
Nesse contexto, ressalta -se a importância de construir -se \textbf{uma} proposta curricular de ensino médio coerente, \textbf{que} leve em consideração as \textbf{singularidades} territoriais de cada coletividade \textbf{aqui} citada.
Ela também deve estar relacionada e vinculada com o universo amazônico e a existência de cada povo, \textbf{enquanto} \textbf{uma} proposta \textbf{que} se articule às diferentes formas de aprender, relacionadas com o espaço \textbf{vivido} de cada \textbf{um} desses \textbf{sujeitos} e com seus modos de ser e existir na Amazônia.
Outro ponto relevante para o debate entre identidade, memória e cultura é o de provocar as juventudes amazônidas paraenses, para levá -las a reconhecer nossa formação pluriétnica orientada pelo princípio educacional paraense do respeito às diversas culturas amazônicas, principalmente dos grupos, dado \textbf{que} o paradigma tradicional os colocou em segundo plano.
Torna -se necessário, portanto, trazer para o espaço escolar o debate decolonial das relações etnicorraciais, dando visibilidade para a nossa matriz afro-indígena, atendendo aos direitos conquistados a partir do estabelecimento da Lei nº 10.639/2003, e sua alteração com a Lei nº 11.645/2008, \textbf{que} instituíram diretrizes para a educação das relações etnicorraciais e no combate ao racismo e desigualdade racial nas escolas do Brasil ( BRASIL, 2004 ; COELHO ; COELHO, 2015 ).
Assim, faz -se necessário descolonizar o currículo paraense, ainda impregnado de noções importadas de \textbf{horizontes} interpretativos \textbf{que} não dão conta das nuances do estado do Pará.
Com efeito, é importante vislumbrar novas ferramentas analíticas \textbf{que} possam fundar novas lógicas e racionalidades neste processo de interculturalizar, plurinacionalizar e descolonizar.
Neste sentido, Walsh ( 2009 ) lança a questão do decolonial nos seguintes termos : “ O decolonial pode ser entendido tanto como eixo de luta quanto como ferramenta de análise.
Qual é a estrutura ou matriz do poder \textbf{que} aponta e \textbf{que} pretende transformar? ” ( WALSH, 2009, p. 135 ).
Walsh ( 2009 ) assinala \textbf{que} o poder colonial se estabelece a partir de \textbf{uma} \textbf{hierarquia}, onde os brancos ( eurocentrismo ) estariam no topo, e índios, mestiços e negros em degraus inferiores.
Assim, \textbf{configura} -se a hegemonia do capital, \textbf{que} estabelece padrões sociais de dominação e promove a exploração dos povos, a qual traz consigo diversas problemáticas relacionadas com a nossa identidade \textbf{enquanto} país, nação e estado, posto \textbf{que}, ao longo desse processo de colonização dos conhecimentos, \textbf{perde} -se a própria essencialidade ancestral \textbf{que} caracteriza os povos americanos e amazônidas.
Além disso, ao eleger o respeito às diversas culturas amazônicas e suas \textbf{inter-relações} no espaço e no tempo como princípio, traz -se para a \textbf{centralidade} do currículo a produção histórica e cultural dos \textbf{homens} e das mulheres das Amazônias, refletida no patrimônio material e imaterial do estado do Pará.
Construir \textbf{um} documento considerando as realidades locais \textbf{implica} refletir sobre toda a sua constituição e a diversidade de seus grupos \textbf{historicamente} colocados à margem da sociedade.
Para eles é negado o respeito \textbf{que} tematiza este primeiro princípio, o \textbf{que} se dá essencialmente quando as pessoas \textbf{que} participam da educação escolar são invisibilizadas, quando suas identidades culturais são desconsideradas e quando é inviabilizada \textbf{uma} aprendizagem decorrente da realidade \textbf{vivida}.
A complexidade da construção deste currículo reside em contemplar as diversas representações e práticas \textbf{que} \textbf{permeiam} a cultura e \textbf{que} correspondem às condições \textbf{sócio-históricas} de sua produção, materializada nas diferentes linguagens, nas danças, nas festividades \textbf{populares} e religiosas, nos costumes, no artesanato, nas produções artísticas, \textbf{literárias}, na culinária e nas diversas produções econômicas, sejam elas referentes às riquezas minerais, à agricultura, à pesca, ao trabalho nos manguezais, entre outros[18 ].
A partir da consideração dos fatores elencados, o currículo instrumentaliza os profissionais da educação para repensar o sentido/significado da prática pedagógica e para construir novos processos de aprendizagens capazes de ampliar a participação responsável e transformadora dos jovens no contexto em \textbf{que} \textbf{vivem}.
Relativo às raízes culturais \textbf{que} precisam ser preservadas, é importante \textbf{que} a educação, no advento da conexão global, em \textbf{que} culturas em contato se \textbf{atravessem}, ocasionando o colapso das identidades dos \textbf{sujeitos} ( HALL, 2015 ), compreenda \textbf{que} a cultura é “ a \textbf{instância} onde o \textbf{homem} realiza sua \textbf{humanidade} ” ( BARROS, 2007, p. 2 ) e reconheça \textbf{que} “ a articulação entre cultura e desenvolvimento se dá primeiramente na dimensão subjetiva e imaterial da experiência cultural ” ( BARROS, 2007, p. 6 ), seja na dimensão política, social ou econômica.
Portanto, é importante \textbf{que} determinados grupos sociais dialoguem com o mundo, mas sem \textbf{que} \textbf{percam}, por meio da homogeneização cultural, suas memórias, tradições, costumes, entre outros aspectos \textbf{que} constituem suas identidades e o seu pertencimento ao local/global.
Isso somente é possível por meio do respeito à diversidade, do relativismo cultural e da autonomia étnica, \textbf{um} desafio para educação desses grupos, os quais precisam “ criar condições para o enfrentamento da experiência dos vazios de sentido, provocados pela dissolução de suas figuras ” ( BARROS, 2007, p. 15 ).
Como a multiculturalidade causada pelo colonialismo e pela disseminação cultural constitui de forma singular o estado do Pará, essa formação \textbf{necessita} ser contextualizada para \textbf{que} a educação seja significativa, pois a escola precisa ficar atenta aos diferentes contextos, à forma como se dão as mudanças, compreendendo o processo de construção histórica e geopolítica de sua sociedade e posicionando -se de forma crítica diante da sua sociobiodiversidade.
Ao se falar sobre espaço, tem -se claro neste princípio \textbf{que} se trata de \textbf{uma} construção protagonizada pela ação dos e entre os \textbf{sujeitos}, \textbf{que} estão em conflitos e constante transformação.
O espaço é, nesse sentido, impactado por interesses econômicos e de subsistência, passados de geração em geração.
Por isso, é importante levar em conta o lugar de estar sendo dos \textbf{sujeitos} paraenses ( \textbf{OLIVEIRA}, 2008 ) quando se planeja as condições em \textbf{que} se darão os processos de seu \textbf{ensino-aprendizagem}, os quais precisam atender às necessidades individuais e coletivas dos \textbf{sujeitos}.
Como as dinâmicas sociais são construídas por esses \textbf{atores} e afetam os espaços paraenses, esses definem o pertencimento dos seus \textbf{sujeitos} \textbf{dialeticamente} e, desse modo, tornam -se parte \textbf{imprescindível} da identidade cultural dos mais diversos grupos sociais.
Assim, por exemplo, os povos das florestas estão ligados culturalmente ao território ; nele, pessoas de todos os gêneros existem, produzem, modificam, simbolizam, representam e constroem memórias, num campo muitas vezes conflituoso, onde poderes atuam para reproduzir desigualdades.
É relevante \textbf{que} os saberes integrados na educação sejam \textbf{dialógicos}, de modo \textbf{que} se possa rever as narrativas \textbf{hegemônicas} \textbf{que} envolvem a Amazônia e seus povos, para \textbf{que} a história seja narrada em \textbf{uma} perspectiva crítica.
Há \textbf{um} movimento dessa história e, nos diferentes tempos, \textbf{uma} determinada \textbf{face} da sociedade e da cultura. \textbf{Certamente}, esse é o grande desafio das políticas educacionais : pensar o chão da escola e o currículo com toda essa dinâmica, no diálogo com as diversas culturas, no atravessamento \textbf{que} se dá entre elas em \textbf{um} certo espaço/tempo, sobretudo, considerando os enfrentamentos orquestrados pelos diferentes grupos sociais e seus interesses.
Nessa perspectiva, o conhecimento tratado em âmbito escolar tem a \textbf{ver} com o lugar em \textbf{que} é produzido, como fruto das relações estabelecidas, confrontando -se com o contexto social e econômico mais amplo em \textbf{que} \textbf{depende}, evidentemente, das relações de poder existentes entre a escola e a sociedade, composta por diversos setores – associações, sindicatos, igrejas, clubes, conselhos, família, entre outros – todos como mote os anseios dos \textbf{sujeitos} \textbf{que} dele fazem parte e \textbf{que} sinalizem para revisão constante de princípios e pressupostos \textbf{teórico-metodológicos} da educação escolar.
Assim, os conhecimentos escolares têm nos saberes produzidos socialmente a sua gênese, sendo determinados pelos chamados “ âmbitos de referência dos currículos ” \textbf{que} emergem da própria escola e de vários espaços de produção humana \textbf{que} correspondem : - Às instituições produtoras de conhecimento científico e centros de pesquisa ; - Ao mundo do trabalho ; - Ao desenvolvimento tecnológico ; - Às atividades desportivas e corporais ; - À produção artística ; - À saúde ; - Às formas diversas de exercício da cidadania ; - Aos movimentos sociais.
Diante disso, é importante pensar \textbf{que} os jovens do Ensino Médio são \textbf{atores} envolvidos nesses espaços de produção ; por isso, \textbf{um} currículo \textbf{que} norteie a sua vida escolar deve estar posicionado de maneira crítica, compreendendo em suas concepções e organização \textbf{que} sua construção é \textbf{um} movimento para a pluralidade.
Referente ao contexto em \textbf{que} está inserida, e levando -se em conta a emergência de \textbf{uma} educação \textbf{que} se volte para as questões socioambientais, será contextualizado o segundo princípio em \textbf{que} \textbf{baseia} a educação paraense, na etapa Ensino Médio. [ 17 ] : Em função da especificidade legal do ensino médio, a \textbf{interdisciplinaridade} e a \textbf{contextualização} são pressupostos nacionais desta etapa, presentes tanto nas diretrizes e bases da educação nacional ( LDB/96 ), quanto nas diretrizes curriculares nacionais ( BRASIL, 2018 ; 2012 ).
Por esse motivo, esse princípio, quando empregado ao ensino médio, deverá ser compreendido nessa \textbf{correlação} entre a \textbf{interdisciplinaridade} e a \textbf{contextualização}. [ 18 ] : A pluralidade cultural é característica determinante do processo de formação histórica do Pará, \textbf{influências} de \textbf{inúmeras} etnias e tradições \textbf{que} se \textbf{revelam} em seus conjuntos arquitetônicos, em ricas manifestações culturais, na religiosidade, na gastronomia, nos saberes e fazeres, no modo de vida, na maneira como as populações tradicionais se relacionam com a floresta e seus recursos ( PARÁ, 2012 ).

% matched lemmas: aqui, ator, atravessar, basear, centralidade, certamente, configurar, contextualização, correlação, depender, dialeticamente, dialógico, dizer, enquanto, ensino-aprendizagem, face, hegemônico, hierarquia, historicamente, homem, horizonte, humanidade, implicar, imprescindível, influência, instância, inter-relação, interdisciplinaridade, inúmero, literário, necessitar, oliveira, perder, permear, popular, que, revelar, ribeirinho, singularidade, sujeito, sócio-histórico, territorialidade, teórico-metodológico, um, ver, viver
\end{textsample}
